\chapter{Integral compactness and mod-two compatibility}\label{sec:integral-varifold-compactness}
Weak convergence of measures alone can lose the integer structure of a surface. First-variation bounds provide an additional geometric compactness principle. Compare the current compactness theorem in the currents chapter with the following varifold result.

\begin{theorem}[Allard compactness for integral varifolds]\label{thm:allard-compactness}
Let $U\subset\mathbb R^N$ be open, and let $V_j$ be integral $m$-varifolds in $U$ with locally bounded first variation. Assume that for every relatively compact open $W\Subset U$,
\[
 \sup_j\bigl(\|V_j\|(W)+\|\delta V_j\|(W)\bigr)<\infty.
\]
Here $\|\delta V_j\|$ is the total variation measure of the first-variation functional. A subsequence converges as varifolds to an integral $m$-varifold $V$ with locally bounded first variation. Integrality of the limit is the substantive additional conclusion beyond weak compactness of Radon measures. Source: \cite{Allard72}, Theorem 6.4; see also \cite{White09}, Section 3. The proof remains external.
\end{theorem}

\begin{definition}[The mod-two chain of an integral varifold]\label{def:varifold-mod-two}
For an integral varifold $V=V(S,\theta)$, let $[V]$ denote the rectifiable flat chain with coefficients in $\mathbb Z_2$ obtained by reducing the integer multiplicity $\theta$ modulo two. No orientation is required. Thus its density is $\theta\bmod 2$ at almost every point of $S$. Varifold convergence concerns measures on the Grassmann bundle, whereas flat convergence concerns chains; neither convergence is identified with the other by this definition. See \cite{White09}, Section 1.
\end{definition}

\begin{theorem}[Compatibility of varifold and mod-two limits]\label{thm:white-mod-two}
Suppose integral $m$-varifolds $V_j$ in an open $U\subset\mathbb R^N$ converge to $V$ as varifolds, with uniformly locally bounded first variations. If $\partial[V_j]$ converges in the local flat topology, then $[V_j]$ converges in that topology to $[V]$. In particular, the integral limit is provided by Theorem \ref{thm:allard-compactness}, and the flat-chain limit agrees with its multiplicity modulo two. Source: \cite{White09}, Theorems 1.1 and 3.3, using the local conventions in Section 4. The boundary convergence assumption is retained. The proof remains external.
\end{theorem}

\begin{example}[Even multiplicity is invisible modulo two]\label{ex:even-multiplicity}
If $P$ is an $m$-plane, the integral varifold $2|P|$ has twice the weight of $|P|$, but $[2|P|]=0$. Consequently the mod-two chain alone does not determine the area or multiplicity of the varifold. This elementary example explains why the two forms of convergence in Theorem \ref{thm:white-mod-two} carry different information; see \cite{White09}, Section 2, for examples involving sequences.
\end{example}

\chapter{Compactness with bounded area and index}\label{sec:bounded-index-compactness}
The Morse index measures the number of independent unstable directions, not the size of the second fundamental form. Its local consequence is particularly useful in compactness arguments.

\begin{lemma}[Disjoint unstable regions consume index]\label{lem:disjoint-unstable-regions}
Let $\Sigma$ be a smooth closed two-sided minimal hypersurface with index at most $I$. Among any $I+1$ open subsets of $\Sigma$ with pairwise disjoint closures, at least one is stable for normal variations compactly supported in that subset.
\end{lemma}
\begin{proof}
Otherwise choose nonzero smooth compactly supported functions $f_1,\ldots,f_{I+1}$, one in each subset, with $Q(f_i,f_i)<0$. The disjoint supports make every cross term $Q(f_i,f_j)$ vanish for $i\ne j$. Hence
\[
 Q\left(\sum_i a_if_i,\sum_i a_if_i\right)=\sum_i a_i^2Q(f_i,f_i)<0
\]
for every nonzero coefficient vector. The functions are linearly independent, producing an $(I+1)$-dimensional negative subspace, contrary to the index bound. This is the localization mechanism used in \cite{Sharp17}, Lemma 3.1.
\end{proof}

\begin{theorem}[Sharp compactness]\label{thm:sharp-compactness}
Let $(M^{n+1},g)$ be smooth and closed, with $2\leq n\leq6$. Let $\Sigma_j$ be smooth, connected, closed, embedded minimal hypersurfaces satisfying
\[
 \operatorname{Area}(\Sigma_j)\leq\Lambda,\qquad
 \operatorname{index}(\Sigma_j)\leq I
\]
for fixed finite $\Lambda$ and integer $I\geq0$. After a subsequence, $|\Sigma_j|$ converges to $m|\Sigma|$ for a positive integer $m$ and a smooth closed embedded minimal hypersurface $\Sigma$, with area at most $\Lambda$ and index at most $I$. The convergence is smooth and graphical away from a set $Y\subset\Sigma$ of at most $I$ points. If $m=1$, it is smooth and graphical everywhere. If $m\geq2$ and $\Sigma$ is two-sided, $\Sigma$ is stable. Source: \cite{Sharp17}, Theorem 2.3, pp.~321--322. The more detailed nullity conclusions are not needed here. The proof remains external.
\end{theorem}

\begin{remark}[Where regularity enters]\label{rem:sharp-regularity-input}
Lemma \ref{lem:disjoint-unstable-regions} limits the number of regions where stability can fail. It does not itself give a curvature estimate. Theorem \ref{thm:stable-curvature-bound} now states the smooth Euclidean curvature estimate explicitly. Theorem \ref{thm:separated-sheets-jacobi} and Proposition \ref{prop:positive-jacobi-stability} explain how separated sheets imply stability of a limit. Sharp's argument uses the local stable-hypersurface estimates of Schoen--Simon \cite{SS81}; the initial use is stated in \cite{Sharp17}, Theorem 2.1. For singular stable hypersurfaces, stability must be accompanied by the appropriate assumptions on the singular set. Theorem \ref{thm:removable-stability-capacity} now proves stability across finite punctures on an already smooth hypersurface, and Theorem \ref{thm:index-smooth-varifold-limit} proves index persistence when the weak limit is already smooth and two-sided. These deductions do not supply the regularity of a singular support. Neither stability nor a bound on index alone asserts that an arbitrary stationary varifold is smooth.
\end{remark}
